\chapter*{第 4 章 流体动力学基本方程}
\addcontentsline{toc}{chapter}{第 4 章 流体动力学基本方程}
\markboth{第 4 章 流体动力学基本方程}{}

\begin{tishi}
\textbf{本章是全卷计算题的头号重点}（E1 slide33"第四章 流体动力学基础＝考试计算题重点"，E2"计算题占 50\% 分值"）。
伯努利方程与动量方程几乎年年必考，务必做到\textbf{能独立写出选断面、基准面、压强基准与适用条件}。
做题要求：选择题也要\textbf{写出判断依据}；伯努利方程类题必须写明\textbf{选断面（为什么）、基准面、压强基准（绝对/相对）、
适用条件（不可压缩、恒定流、两点取在渐变流断面上）}；动量方程类题必须写明\textbf{控制体、受力分析、投影方向的正负号约定}，
最后明确回答"是流体对壁面还是壁面对流体"。全部题目都要有实质解答，不得写"略"。
教材原书本章习题共 16 题（4.1--4.16），题号下的图按教材原图，本册用〔图示〕文字补充已知量，\textbf{绝不编造尺寸}。
\end{tishi}

\begin{ti}
\sloppy
\emergencystretch=2em

\item 选择题。（1）如习题 4.1（1）图所示，等直径管，考虑损失，$A$-$A$ 断面为过流断面，$B$-$B$ 断面为水平面，试确定 1、2、3、4 各点的物理量有以下关系（\hspace{1cm}）。
A. $z_1+\dfrac{p_1}{\rho g}=z_2+\dfrac{p_2}{\rho g}$\quad B. $p_3=p_4$\quad C. $p_1=p_2$\quad D. $z_3+\dfrac{p_3}{\rho g}=z_4+\dfrac{p_4}{\rho g}$
（2）伯努利方程中 $z+\dfrac{p}{\rho g}+\dfrac{\alpha v^2}{2g}$ 表示（\hspace{1cm}）。
A. 单位质量流体具有的机械能\quad B. 单位重量流体具有的机械能\quad C. 单位体积流体具有的机械能\quad D. 通过过流断面流体的总机械能
（3）水流的运动方向是（\hspace{1cm}）。
A. 从高处往低处流\quad B. 从压强高处往压强低处流\quad C. 从流速大的地方向流速小的地方流\quad D. 从单位重量流体机械能高的地方向低的地方流
（4）如习题 4.1（4）图所示，水平放置的渐扩管，如忽略水头损失，断面形心点的压强有以下关系（\hspace{1cm}）。
A. $p_1>p_2$\quad B. $p_3=p_4$\quad C. $p_1<p_2$\quad D. 不定
\tuyi{习题 4.1（1）图：{\bfseries 等直径圆管}内水流自左向右流动；{\bfseries $A$-$A$ 断面为过流断面}（与管轴垂直的竖直线），其上标出点 1（在管轴上方）与点 2（在管轴下方）；{\bfseries $B$-$B$ 断面为水平面}（沿管轴的水平线），其上标出点 3（在 $A$-$A$ 左侧管轴上）与点 4（在 $A$-$A$ 右侧管轴上）；1、2 两点在 $A$-$A$ 上，3、4 两点在 $B$-$B$ 上。\par
习题 4.1（4）图：水平放置的{\bfseries 渐扩管}，小端（左）为断面 1-1，形心压力 $p_1$；大端（右）为断面 2-2，形心压力 $p_2$；管内水流自左向右，管中路线上标注流向箭头。}
\xstarfull{4}\yiju{E4 2022 年 818 单选第 6 题"伯努利方程所选计算点位置"（答案：应取在渐变流断面上且使未知量最少处）；E1 slide33"第四章＝考试计算题重点"；教材式 4.17、4.18 与"测压管水头线"（E6）；本章选择题未在真题中整章出现过，故压至四星}\nandu{中}

\item 利用毕托管原理测量输水管中的流量。已知输水管直径 $d=200\ \mathrm{mm}$，水银压差计读数 $\Delta h=60\ \mathrm{mm}$，如习题 4.2 图所示，若输水管断面平均流速 $v=0.84u_A$，式中 $u_A$ 是管轴上未受扰动的 $A$ 点的流速。试确定输水管中的流量 $Q$。
\tuyi{习题 4.2 图：一段输水管，泵向流动方向由{\bfseries $u_A$} 与箭头标出；管轴上取点 $A$，由 $A$ 点引出{\bfseries 毕托管}的两根测压细管（一根测静压、一根测总压），两管接至一个{\bfseries U 形水银压差计}，压差计内装有水银（图中以斜线阴影表示），两臂水银面的{\bfseries 高差为 $\Delta h=60\ \mathrm{mm}$}；{\bfseries 输水管直径 $d=200\ \mathrm{mm}$}（图中在管下方以 $d$ 标注）。}
\xstarfull{5}\yiju{E4 2015 年计算题第 5 题"文丘里流量计"（同属毕托管类测流原理，教材式 4.11、4.12、4.21）；E2"第四章主要考点需掌握计算套路"；E6 教材 4.3 节毕托管与文丘里管例题（式 4.11、4.12）；教材式 4.21 $Q=K\sqrt{\Delta h}$ 即本题的直接依据}\nandu{易}

\item 如习题 4.3 图所示，大水箱中的水经水箱底部的竖管流入大气，竖管直径为 $d_1=200\ \mathrm{mm}$，管道出口处为收缩喷嘴，其直径 $d_2=100\ \mathrm{mm}$，不计水头损失，求管道的泄流量 $Q$ 及 $A$ 点相对压强 $p_A$。
\tuyi{习题 4.3 图：一个{\bfseries 大水箱}，箱中水面高于箱底；水由箱底中心的{\bfseries 竖直管}向下流出（流入大气）；竖直管{\bfseries 直径 $d_1=200\ \mathrm{mm}$}（图中在 $A$ 点所在断面标注）；{\bfseries $A$ 点位于竖管上、距水箱内底（管道进口）$4\ \mathrm{m}$}；管道出口处接一段{\bfseries 收缩喷嘴，出口直径 $d_2=100\ \mathrm{mm}$}；{\bfseries 水箱内水面至喷嘴出口所在水平面的高度为 $7\ \mathrm{m}$}（图中以 7 m 的尺寸线自水面标到喷嘴出口所在的水平面）。}
\xstarfull{5}\yiju{E1 slide33"第四章＝考试计算题重点"；E2"第四章主要考点需掌握计算套路"；E6 教材式 4.17、4.18 与例 4.1（水池有压管道泄水）的同类方法；本题为"箱内竖管＋收缩喷嘴"标准题型，模拟题与题库中反复出现}\nandu{中}

\item 离心式通风机由吸气管吸入空气。如习题 4.4 图所示，吸气管圆筒形部分的直径 $d=200\ \mathrm{mm}$，在这个圆筒形壁上安装一个盛水的测压装置，现量得其中的水面高差 $\Delta h=0.25\ \mathrm{m}$，空气的重度为 $12.6\ \mathrm{N/m^3}$，问此风机在 $1\ \mathrm{s}$ 内的吸气量 $Q$ 为多少？（不计损失）
\tuyi{习题 4.4 图：一段{\bfseries 吸气管}，管右端标"接风机"（气流自左向右被吸入）；管筒形部分的{\bfseries 直径 $d=200\ \mathrm{mm}$}（标注在管体下方）；圆筒壁上开一测压孔并接出一根{\bfseries U 形测压管}（内盛水），测压管一端接管壁、另一端通大气，测得{\bfseries 管内水面高差 $\Delta h=0.25\ \mathrm{m}$}（管中水面低于大气端水面）。}
\xstarfull{4}\yiju{E5 题库与 2015 年试题中"用测压管与毕托管原理测流速"反复出现；E1 slide33"第四章＝考试计算题重点"；E6 教材式 4.11、4.12 与毕托管测量原理（气体用重度给出的口径见教材 1.2 节）；本题为毕托管原理用于气体的标准题型，故四星}\nandu{中}

\item 如习题 4.5 图所示，虹吸管从水池引水至 $C$ 端流入大气，已知 $a=1.6\ \mathrm{m}$，$b=3.6\ \mathrm{m}$。若不计损失，试求：（1）管中流速 $v$ 及 $B$ 点的绝对压强 $p_B$；（2）若 $B$ 点绝对压强水头下降到 $0.24\ \mathrm{m}$ 以下时，将发生汽化，设 $C$ 端保持不动，问欲不发生汽化，$a$ 不能超过多少？
\tuyi{习题 4.5 图：一个{\bfseries 水池}（图中右侧以斜线阴影示意池壁），左上方为池中水面；一根{\bfseries 虹吸管}自池中伸入水下，向上翻过池壁、在池壁外侧达到{\bfseries 最高点 $B$}，再向下倾斜，末端为{\bfseries 出口 $C$}（流入大气，出口标有表示自由出流的短斜线）；{\bfseries $B$ 点高出池中水面的高度为 $a=1.6\ \mathrm{m}$}（图中以 $a$ 标注在 $B$ 点与水面标高线之间）；{\bfseries $C$ 端低于池中水面的深度为 $b=3.6\ \mathrm{m}$}（图中以 $b$ 标注在 $C$ 点与水面标高线之间）。}
\xstarfull{5}\yiju{E4 2022 年 818 单选第 7 题"虹吸管最高处压强与大气压的关系"（答案：最高处绝对压强小于大气压）；E4 2022 年单选第 8 题圆柱形直角外管嘴正常工作的临界高度（同属"虹吸/管嘴的负压与汽化"考点）；E5 题库选择"虹吸管最高处压强小于大气压"；E1 slide33"第四章＝考试计算题重点"}\nandu{难}

\item 如习题 4.6 图所示，用抽水量 $Q=24\ \mathrm{m^3/h}$ 的离心水泵由水池抽水，水泵的安装高程 $h_s=6\ \mathrm{m}$，吸水管的直径为 $d=100\ \mathrm{mm}$，如水流通过进口底阀、吸水管路、$90^{\circ}$ 弯头至泵叶轮进口的总水头损失为 $h_w=0.4\ \mathrm{mH_2O}$，求该泵叶轮进口处的真空度 $p_v$。
\tuyi{习题 4.6 图：右侧为一个{\bfseries 水池}（图中以阴影线示意池底与池壁），池中水面低于左侧的泵房地面；吸水管自池中（管口装有带孔的{\bfseries 底阀}）向上穿过池壁，经一个 $90^{\circ}$ 弯头接至{\bfseries 离心水泵}（图中以圆形表示泵体）；泵的出水管向右引出并在管路上标出{\bfseries 流量 $Q$}；{\bfseries 水泵的安装高程（泵轴线高于水池液面）$h_s=6\ \mathrm{m}$}（图中以 $h_s$ 标注在泵轴线与池中水面标高线之间）；{\bfseries 吸水管直径 $d=100\ \mathrm{mm}$}。}
\xstarfull{5}\yiju{E1 slide33"第四章＝考试计算题重点"；E2"第四章主要考点需掌握计算套路"；E6 教材式 4.22、4.23（有能量输入的总流伯努利方程）；本题为"水泵吸入口真空度（允许吸上真空高度）"标准题型，E5 题库选择考过"泵允许吸上真空高度 $H_s$"，故五星}\nandu{中}

\item 如习题 4.7 图所示，高压水箱的泄水管，当阀门关闭时，测得安装在此管路上的压力表读数为 $p_1=280\ \mathrm{kPa}$，当阀门开启后，压力表上的读数变为 $p_2=60\ \mathrm{kPa}$，已知此泄水管的直径 $D=25\ \mathrm{mm}$，求每小时的泄水流量。（不计水头损失）
\tuyi{习题 4.7 图：一个{\bfseries 高压水箱}（图中以箱体示意，箱中水面高于泄水管），箱底接出一段{\bfseries 水平泄水管}向右侧流出（出口为自由出流，流入大气，图中出口处画有向下翻折的短管与出流箭头）；水平管路上装有一个{\bfseries 压力表}（图中以带指针的圆形表示），表下以 $h$ 标注测点位置；{\bfseries 泄水管直径 $D=25\ \mathrm{mm}$}；阀门画在管路下游（压力表与出口之间）。}
\xstarfull{5}\yiju{E1 slide33"第四章＝考试计算题重点"；E2"第四章主要考点需掌握计算套路"；E6 教材式 4.17、4.18 与"总水头线、测压管水头线"的物理含义：\textbf{阀门关闭时管内流速为零，压力表读数即为该断面的静压}；本题为"启阀前后压力表读数差转化为动能"标准题型，历年模拟卷常见}\nandu{中}

\item 如习题 4.8 图所示，由水池通过等直径虹吸管输水，$A$ 点为虹吸管进口处，$H_A=0$；$B$ 点为虹吸管中与水池液面齐高的部位，$H_B=6\ \mathrm{m}$；$C$ 点为虹吸管中的最高点，$H_C=7\ \mathrm{m}$；$D$ 点为虹吸管的出口处，$H_D=4\ \mathrm{m}$。若不计流动中的能量损失，求虹吸管的断面平均流速和 $A$、$B$、$C$ 各断面上的绝对压强。
\tuyi{习题 4.8 图：右侧为{\bfseries 水池}，池中水面以标高线标出；一根{\bfseries 等直径虹吸管}自池中伸入水下（进口断面为 $A$，标高 $H_A=0$），向上翻越池壁后在池外形成{\bfseries 最高点 $C$}（标高 $H_C=7\ \mathrm{m}$），再向下弯折（管上标出 $B$ 点，标高 $H_B=6\ \mathrm{m}$，位于与水池液面齐高处），末端为{\bfseries 出口 $D$}（标高 $H_D=4\ \mathrm{m}$，图中 $D$ 端低于池中液面）；图中以竖直尺寸线自池中水面标高线分别标到 $A$、$B$、$C$、$D$ 四点。}
\xstarfull{5}\yiju{E4 2022 年 818 单选第 7 题"虹吸管最高处压强与大气压的关系"；E5 题库选择"虹吸管最高处压强小于大气压"；E1 slide33"第四章＝考试计算题重点"；E6 教材式 4.18 与例 4.1 的同类方法；本题是"虹吸管逐点求绝对压强"的典型综合题}\nandu{难}

\item 如习题 4.9 图所示，水泵的提水高度 $z=20\ \mathrm{m}$，抽水流量 $Q=35\ \mathrm{L/s}$，已知吸水管和压水管的直径相同，$d=180\ \mathrm{mm}$，离心泵的效率 $\eta_1=0.82$，电动机的效率 $\eta_2=0.95$，设总水头损失 $h_w=1.5\ \mathrm{mH_2O}$，求电动机应有的功率 $P$。
\tuyi{习题 4.9 图：左侧为一个{\bfseries 水池}（图中以阴影线示意池中水面与池底），一根{\bfseries 吸水管}自池中伸入水下（管上标断面 1-1），向上接至{\bfseries 离心水泵}（图中以圆形表示泵体，泵上方标"水泵"）；泵的{\bfseries 压水管}向右上倾斜引出（管上标断面 2-2），末端为压力出流；图中以竖直尺寸线标出{\bfseries 提水高度 $z=20\ \mathrm{m}$}（自池中水面标高线到压水管出口断面 2-2 所在的高度）。}
\xstarfull{5}\yiju{E1 slide33"第四章＝考试计算题重点"；E2"计算题占 50\% 分值、第四章主要考点需掌握计算套路"；E6 教材式 4.23、4.24、4.25 与例 4.2（数据完全相同：$z=20\ \mathrm{m}$、$Q=35\ \mathrm{L/s}$、$\eta_1=0.82$、$\eta_2=0.95$、$h_w=1.5\ \mathrm{mH_2O}$）；2026 回忆版第 5 题即教材例题改编，故教材例题类必练}\nandu{中}

\item 闸下出流，平板闸门宽 $b=2\ \mathrm{m}$，闸前水深 $h_1=4\ \mathrm{m}$，闸后水深 $h_2=0.5\ \mathrm{m}$，如习题 4.10 图所示，出流量 $Q=8\ \mathrm{m^3/s}$，不计摩擦阻力，试求水流对闸门的作用力 $F$。
\tuyi{习题 4.10 图：左侧为{\bfseries 上游渠道}，水深 $h_1=4\ \mathrm{m}$（图中以竖直尺寸线自渠底标到自由水面标高线）；渠道中部为{\bfseries 平板闸门}（图中以带阴影的竖直线示意，闸门底部提起一段距离形成闸下孔口）；闸门下游水深 $h_2=0.5\ \mathrm{m}$（图中以竖直尺寸线自渠底标到下游水面）；渠底为水平面；水流自左向右经闸下射出；{\bfseries 闸门宽 $b=2\ \mathrm{m}$（垂直于纸面）}。}
\xstarfull{5}\yiju{E1 slide33"第五章＝考试计算题重点、第四章＝考试计算题重点"；E2"计算题占 50\% 分值"；E6 教材式 4.30、4.31（总流动量方程与其分量式）；本题为"闸下出流动量"标准题型，与 4.11 溢流坝推力同源，属动量方程必练题}\nandu{难}

\item 溢流坝宽度为 $B$（垂直于纸面），上游和下游水深分别为 $h_1$ 和 $h_2$，如习题 4.11 图所示，不计水头损失，试推证坝体受到的水平推力 $F=\dfrac{\rho gB(h_1-h_2)^3}{2(h_1+h_2)}$。
\tuyi{习题 4.11 图：左侧为{\bfseries 上游}，水深 $h_1$（图中以竖直尺寸线自坝基高程标到上游水面），水流以 $v_1$ 自左向右接近坝体（图中以水平箭头标注）；渠道中部为{\bfseries 溢流坝}（图中以带阴影的坝体示意）；坝下游水深 $h_2$（图中以竖直尺寸线自坝基高程标到下游水面），水以 $v_2$ 向右下方流去（图中以斜箭头标注）；断面 1-1 取在坝上游、断面 2-2 取在坝下游；{\bfseries 坝宽 $B$（垂直于纸面）}。}
\xstarfull{4}\yiju{E1 slide33"第四章＝考试计算题重点"；E2"第四章主要考点需掌握计算套路"；E6 教材式 4.30、4.31；E4 历年真题中未见直接考"溢流坝推力推证"，但"动量方程＋连续性方程＋能量方程联立"的方法与 4.10 完全同源，故四星}\nandu{难}

\item 水流经水平弯管流入大气，已知 $d_1=100\ \mathrm{mm}$，$d_2=75\ \mathrm{mm}$，$v_1=1.5\ \mathrm{m/s}$，$\theta=30^{\circ}$，如习题 4.12 图所示。若不计水头损失，试求水流对弯管的作用力 $F_x$、$F_y$。
\tuyi{习题 4.12 图：一段{\bfseries 水平放置的收缩弯管}（管径由粗变细）；进口断面为 1-1，{\bfseries 进口直径 $d_1=100\ \mathrm{mm}$}，水流以{\bfseries 速度 $v_1=1.5\ \mathrm{m/s}$} 沿管轴进入（图中以水平箭头标注）；出口断面为 2-2，{\bfseries 出口直径 $d_2=75\ \mathrm{mm}$}，水流以{\bfseries 速度 $v_2$} 射入大气（图中以向右下方的斜箭头标注）；弯管的{\bfseries 转折角 $\theta=30^{\circ}$}（图中标注在进口轴线与出口轴线的夹角处）。}
\xstarfull{5}\yiju{E4 2015 年计算题第 4 题"水平弯管动量方程求 $F_x$、$F_y$"（与本题同型）；E1 slide33"第四章＝考试计算题重点"；E2"第四章主要考点需掌握计算套路"；E6 教材式 4.30、4.31 与例 4.3（水平变径弯管动量方程）；近年真题重复率很高（E1 slide25），故五星}\nandu{中}

\item 水平分岔管路，$d_1=500\ \mathrm{mm}$，$d_2=400\ \mathrm{mm}$，$d_3=300\ \mathrm{mm}$，$Q_1=0.35\ \mathrm{m^3/s}$，$Q_2=0.2\ \mathrm{m^3/s}$，$Q_3=0.15\ \mathrm{m^3/s}$，表压强 $p_1=8000\ \mathrm{N/m^2}$，夹角 $\alpha=45^{\circ}$，$\beta=30^{\circ}$，如习题 4.13 图所示。忽略水头损失，求水流对分岔管的作用力 $F_x$、$F_y$。
\tuyi{习题 4.13 图：一根水平{\bfseries 主管}（断面 1-1）自左向右进水，{\bfseries $d_1=500\ \mathrm{mm}$、$Q_1=0.35\ \mathrm{m^3/s}$}（断面 1-1 处标出 $Q_1$ 与流向箭头）；主管在{\bfseries 分岔处}分为两支：{\bfseries 支管 2（断面 2-2）与主管轴线成夹角 $\alpha=45^{\circ}$}，{\bfseries $d_2=400\ \mathrm{mm}$、$Q_2=0.2\ \mathrm{m^3/s}$}（向斜上方引出）；{\bfseries 支管 3（断面 3-3）与主管轴线成夹角 $\beta=30^{\circ}$}，{\bfseries $d_3=300\ \mathrm{mm}$、$Q_3=0.15\ \mathrm{m^3/s}$}（向斜下方引出）；两个支管出口均排入大气。}
\xstarfull{5}\yiju{E4 2015 年计算题第 4 题"水平弯管动量方程"与本题同属"多进出口动量方程"考点；E1 slide33"第四章＝考试计算题重点"；E2"第四章主要考点需掌握计算套路"；E6 教材式 4.30、4.31 与多进出口形式（教材 4.5 节对分流、汇流的处理）；近年真题重复率很高}\nandu{难}

\item 如习题 4.14 图所示，流量为 $Q$、平均流速为 $v$ 的射流，冲击直立平板后分成两股，一股沿板面直泻而下，流量为 $Q_1$，另一股以倾角 $\alpha$ 射出，流量为 $Q_2$，两股分流的平均速度均等于 $v$。若不计摩擦力和重力影响，试推证：固定平板所需的外力 $F=\rho Qv\left(1-\sqrt{\dfrac{1-Q_1/Q_2}{1+Q_1/Q_2}}\right)$。
\tuyi{习题 4.14 图：左侧一根水平{\bfseries 射流管}（断面 0-0），水流以{\bfseries 流量 $Q$、平均流速 $v$} 水平向右射出（图中以水平箭头标注）；射流{\bfseries 垂直冲击一块直立平板}（图中以竖直双线示意）；冲击后分流为两股：{\bfseries 一股沿板面直泻而下}（图中以沿板面向下的竖箭头标注，流量 $Q_1$、速度 $v$）；{\bfseries 另一股以倾角 $\alpha$ 斜射出去}（图中以与板面成 $\alpha$ 角的斜箭头标注，流量 $Q_2$、速度 $v$）；两股的速度大小均等于来流速度 $v$。}
\xstarfull{4}\yiju{E6 教材式 4.30、4.31 与例 4.4（水平射流冲击斜置平板：$Q_0=0.06\ \mathrm{m^3/s}$、$F'=0.312\ \mathrm{kN}$，沿板方向流量比按动量方程解出）；E1 slide33"第四章＝考试计算题重点"；E4 历年真题中未见"射流冲击平板"独立成题，故四星}\nandu{难}

\item 已知离心式通风机叶轮的转速 $n=1500\ \mathrm{r/min}$，叶轮进口直径 $d_1=480\ \mathrm{mm}$，进口角 $\beta_1=60^{\circ}$，入口宽度 $b_1=105\ \mathrm{mm}$，出口直径 $d_2=600\ \mathrm{mm}$，出口角 $\beta_2=120^{\circ}$，出口宽度 $b_2=84\ \mathrm{mm}$，流量 $Q=12000\ \mathrm{m^3/h}$，如习题 4.15 图所示。试求：（1）叶轮进出口空气的牵连速度 $u_1$、$u_2$，相对速度 $w_1$、$w_2$，绝对速度 $v_1$、$v_2$；（2）单位重量空气通过叶轮所获得的能量 $H$。
\tuyi{习题 4.15 图：{\bfseries 离心式通风机叶轮}的进出口{\bfseries 速度三角形}；图中左侧为一对进口速度三角形（含进口角 $\beta_1$、相对速度 $w_1$、绝对速度 $v_1$、牵连速度 $u_1$），右侧为一对出口速度三角形（含出口角 $\beta_2$、相对速度 $w_2$、绝对速度 $v_2$、牵连速度 $u_2$）；图中以圆弧与箭头表示叶轮旋转方向（角速度 $\omega$）；进口速度三角形由 $w_1$、$u_1$、$v_1$ 围成、$\beta_1$ 为 $w_1$ 与 $u_1$ 反向延长线之间的夹角，出口速度三角形由 $w_2$、$u_2$、$v_2$ 围成、$\beta_2$ 为 $w_2$ 与 $u_2$ 反向延长线之间的夹角。}
\xstarfull{4}\yiju{E6 教材式 4.32、4.33 与例 4.5（离心风机叶轮：$n=1725\ \mathrm{r/min}$、$d_1=125\ \mathrm{mm}$、$d_2=300\ \mathrm{mm}$、$\alpha_1=90^{\circ}$、$\beta_2=30^{\circ}$、$b_1=b_2=25\ \mathrm{mm}$、$Q=372\ \mathrm{m^3/h}$，答案 $u_1=11.28$、$v_1=10.53\ \mathrm{m/s}$、$\beta_1=43.03^{\circ}$、$u_2=27.08$、$v_{2n}=4.39$、$v_{2u}=19.48$、$v_2=19.97\ \mathrm{m/s}$、$\alpha_2=12.72^{\circ}$、$H_T=53.83\ \mathrm{m}$）——本题与例 4.5 方法完全相同、只换数据；E1 slide33"第四章＝考试计算题重点"；4.4 节动量矩方程本身真题未独立成题，故四星}\nandu{难}

\item 臂长 $l_1=1.2\ \mathrm{m}$，$l_2=1.5\ \mathrm{m}$ 的旋转式洒水器，喷口直径 $d=25\ \mathrm{mm}$，每个喷口的水流量 $Q=3\times10^{-3}\ \mathrm{m^3/s}$，如习题 4.16 图所示。若不计摩擦阻力，试确定洒水器的转速 $\omega$。
\tuyi{习题 4.16 图：一个{\bfseries 旋转式洒水器}（两根水平臂在同一铅直平面内、绕中心铅直轴转动）；两臂长度分别为{\bfseries $l_1=1.2\ \mathrm{m}$（左臂）}与{\bfseries $l_2=1.5\ \mathrm{m}$（右臂）}（图中以 $l_1$、$l_2$ 分别标注自转轴到两臂喷口处的长度）；两臂末端各有{\bfseries 一个喷口，喷口直径 $d=25\ \mathrm{mm}$}；左臂喷口朝下（图中以向下的箭头标注出流）、右臂喷口朝上（图中以向上的箭头标注出流），两股射流方向相反；两臂在转轴处相接，转轴以{\bfseries 角速度 $\omega$} 旋转（图中以弧形箭头标注旋转方向）。}
\xstarfull{4}\yiju{E6 教材式 4.32（动量矩方程）与 4.4 节叶轮机械基本方程；E1 slide33"第四章＝考试计算题重点"；E2"第四章主要考点需掌握计算套路"；E4 历年真题中未见"旋转式洒水器"独立成题，但动量矩方程属教材要求掌握内容，故四星}\nandu{难}

\end{ti}

\begin{tishi}
\textbf{〔图示说明〕}本章原书习题图较多，本册一律用〔图示〕文字精确描述已知量，\textbf{不重绘图、不编造尺寸}。
凡题面尺寸与图不符者，以\textbf{教材正文例题口径}为准，并在解答中写明所取数值。

\textbf{〔存疑处（解题时请以题图为准）〕}
\textbf{第 4.3 题：}竖管上 $A$ 点距箱底 $4\ \mathrm{m}$、液面至喷嘴出口 $7\ \mathrm{m}$ 由习题图尺寸线读出；
若你的题图中 $7\ \mathrm{m}$ 量到的是喷嘴出口以上的其他位置，请按图中实际尺寸线取值，方法与解答完全相同。\\
\textbf{第 4.8 题：}$H_D=4\ \mathrm{m}$ 是题目原文给出的标高值，而习题图表示 $D$ 端\emph{低于}池中液面，
故解答中按"标高"取 $z_D=-4\ \mathrm{m}$（以水池液面为基准面）；若你的题图把 $H_D$ 画在液面之上，请改取 $+4\ \mathrm{m}$，
其余方法与解答不变。\\
\textbf{第 4.5、4.8 题（大气压取值）：}教材原题未给大气压数值，解答中按 $p_a=98\ \mathrm{kPa}$ 计算，
同时给出 $p_a=101.3\ \mathrm{kPa}$ 时的结果；答题时按卷面给定值代入，并写明所取值。\\
\textbf{第 4.15 题：}$\beta_2=120^{\circ}$ 的基准依习题图判定（图中 $\beta$ 量自切向量）。
解答按教材例 4.5 的口径 $v_{2u}=u_2-w_2\cos\beta_2$ 计算；若按"$\beta$ 取切向锐角"的量法，$H_T$ 会不同，
解答中已并列给出两种结果与说明。
\end{tishi}
